\DocumentMetadata{
  pdfversion=2.0,
  pdfstandard=ua-2,
  lang=en-US,
  tagging=on,
  tagging-setup={math/setup=mathml-SE,tabsorder=structure}
}
\documentclass[11pt,letterpaper]{article}
\usepackage[margin=0.68in,includefoot]{geometry}
\usepackage{newcomputermodern}
\usepackage{amsmath,amscd,mathtools}
\usepackage{tikz,adjustbox}
\providecommand{\square}{\mdlgwhtsquare}
\usepackage[hidelinks]{hyperref}
\hypersetup{
  pdftitle={Algebra First-Year Exam, January 2018, Part 1},
  pdfauthor={University of Florida Department of Mathematics},
  pdfsubject={Graduate mathematics examination},
  pdfdisplaydoctitle=true,
  bookmarksopen=true
}
\setcounter{secnumdepth}{0}
\setlength{\parindent}{0pt}
\setlength{\parskip}{0.28em}
\setlength{\emergencystretch}{3em}
\setlength{\leftmargini}{2.15em}
\setlength{\leftmarginii}{2.65em}
\setlength{\leftmarginiii}{3em}
\pagestyle{empty}
\raggedbottom
\allowdisplaybreaks
\NewTaggingSocketPlug{block/list/label}{examproblem}{%
  \tagstructbegin{tag=Lbl}\tagstructend%
  \tagstructbegin{tag=\LBody}%
  \tagstructbegin{tag=\ProblemHeadingTag,title-o={\ProblemHeadingText}}%
  \tagmcbegin{tag=\ProblemHeadingTag,actualtext={\ProblemHeadingText}}%
  #2%
  \tagmcend%
  \tagstructend%
}
\newenvironment{examproblems}[2]{%
  \def\ProblemHeadingTag{#2}%
  \AssignTaggingSocketPlug{block/list/label}{examproblem}%
  \begin{enumerate}%
  \setcounter{enumi}{#1}%
  \renewcommand{\labelenumi}{\textbf{\theenumi.}}%
  \setlength{\topsep}{0.35em}%
  \setlength{\partopsep}{0pt}%
  \setlength{\itemsep}{0.48em}%
  \setlength{\parsep}{0.18em}%
}{\end{enumerate}}
\newenvironment{examsubparts}{%
  \AssignTaggingSocketPlug{block/list/label}{default}%
  \begin{enumerate}%
  \setlength{\topsep}{0.18em}%
  \setlength{\partopsep}{0pt}%
  \setlength{\itemsep}{0.16em}%
  \setlength{\parsep}{0.1em}%
}{\end{enumerate}}
\newenvironment{examinstructions}{%
  \par\begingroup\itshape
}{%
  \par\endgroup\vspace{0.18em}
}
\makeatletter
\renewcommand\subsection{\@startsection{subsection}{2}{0pt}%
  {-1.25ex plus -.25ex minus -.1ex}%
  {0.55ex plus .1ex}%
  {\normalfont\large\bfseries}}
\renewcommand{\@maketitle}{%
  \newpage\null\vskip -1em%
  {\footnotesize\bfseries
    \href{https://www.ufl.edu/}{University of Florida}, \href{https://math.ufl.edu/}{Department of Mathematics}\par}%
  \vskip -0.5em%
  \tagstructbegin{tag=H1,title={Algebra First-Year Exam, January 2018, Part 1}}%
  \tagpdfparaOff%
  \tagmcbegin{tag=H1}%
    {\LARGE\bfseries\href{https://gma.math.ufl.edu/exams/algebra-first-year/}{Algebra First-Year Exam}, January 2018, Part 1\par}%
  \tagmcend%
  \tagpdfparaOn%
  \tagstructend%
  \vskip 0.25em%
  \tagmcbegin{artifact}%
    \hrule height 1pt%
  \tagmcend%
  \vskip 1em%
}
\makeatother
\title{Algebra First-Year Exam, January 2018, Part 1}
\author{}
\date{}
\begin{document}
\maketitle
\thispagestyle{empty}
\begin{examinstructions}
Answer four problems. Write your answers clearly in complete English sentences. You may quote results (within reason) as long as you state them clearly.
\end{examinstructions}
\begin{examproblems}{0}{H2}
\def\ProblemHeadingText{Problem 1.}
\item
Give an example of each, with a few words of explanation.
\begin{examsubparts}
\item[\textnormal{(a)}]
An abelian simple group.
\item[\textnormal{(b)}]
A nonabelian simple group.
\item[\textnormal{(c)}]
A nilpotent group which isn’t a~\(p\)-group.
\item[\textnormal{(d)}]
A finite group which is neither solvable nor simple.
\end{examsubparts}
\def\ProblemHeadingText{Problem 2.}
\item
Let~\(G\) act on the set~\(A\) and let \(a \in A\). Prove that \(|G \cdot a| = |G : G_a|\). (Note that \(G \cdot a\) is the orbit containing~\(a\), while \(G_a\) is the stabilizer of~\(a\) in~\(G\).)
\def\ProblemHeadingText{Problem 3.}
\item
Let~\(G\) be a finite abelian group of order~\(n\) and let~\(k\) be an integer which is relatively prime to~\(n\). Define \(\tau : G \to G\) by \(\tau(g) = g^k\). Prove that \(\tau \in \operatorname{Aut}(G)\).
\def\ProblemHeadingText{Problem 4.}
\item
Let~\(G\) be a finite group and let~\(H\) and~\(K\) be subgroups of~\(G\). Recall that \(HK = \{xy : x \in H, y \in K\}\).
\begin{examsubparts}
\item[\textnormal{(a)}]
Prove the following formula for the cardinality of the set \(HK\):
\[
|HK| = \frac{|H| \cdot |K|}{|H \cap K|}.
\]
\item[\textnormal{(b)}]
Give an example in which \(|HK|\) does not divide \(|G|\).
\end{examsubparts}
\def\ProblemHeadingText{Problem 5.}
\item
Let~\(G\) be a group and let~\(A\) be an abelian group.
\begin{examsubparts}
\item[\textnormal{(a)}]
Define the commutator subgroup \(G' = [G,G]\) of~\(G\).
\item[\textnormal{(b)}]
Prove that if \(\phi : G \to A\) is a homomorphism then there is a unique homomorphism \(\overline{\phi} : G/G' \to A\) such that \(\phi = \overline{\phi} \circ \pi\), where \(\pi : G \to G/G'\) is the projection map.
\end{examsubparts}
\def\ProblemHeadingText{Problem 6.}
\item
Determine all groups of order \(99\) up to isomorphism.
\end{examproblems}
\end{document}
